\begin{frame}{왜 ``가중치가 작다''가 ``모델이 단순하다''인가}
  \begin{block}{핵심 메커니즘}
    특징을 평균 0·표준편차 1로 표준화해 두면, $x_j$가 $\Delta$만큼 변할
    때 예측치는 $|w_j|\,\Delta$ 이상으로 변할 수 있다 — $|w_j|$는 ``특징
    $x_j$가 출력을 흔들어놓을 수 있는 \textbf{강도}''다.
    \vskip0.3em
    모든 $|w_j|$가 작으면 입력이 아무리 요동쳐도 출력은 거의 못 흔들린다.
    4.1에서 ``노이즈 한 점에 예측이 크게 출렁이는 모델'' = 분산이 큰 모델이었는데,
    가중치가 작은 모델은 그 \textbf{정 반대}(흔들림이 억눌림) → 분산이 작다.
  \end{block}
  \begin{itemize}
    \item ``가중치 축소(shrinkage) = 유효 복잡도 감소''가 정규화의 핵심.
    \item 이 관점 때문에 가중치 크기 페널티는 deep learning에서
      \kb{weight decay}(가중치 감쇠)로 그대로 사용된다.
  \end{itemize}
  \vskip0.4em
  \kq{역사: Ridge는 1970년 마쿼트(불량 최소제곱에 이미 1962년부터 사용)와
  호엘·케나드($\lambda I$를 더한 행렬이 손실면에 ``등정(ridge) 모양'' 돌출)의
  이름. Lasso는 26년 뒤, 1996년 티브시라니.}
  {\footnotesize 주의: 이 연결은 \textbf{특징 스케일이 같을 때}만 성립 →
  Ridge/Lasso 전에 \textbf{스케일링을 먼저} 하는 것이 실전 표준.}
\end{frame}
